% Generated draft: requires mathematical and source review.
\documentclass[11pt]{article}
\usepackage[margin=25mm]{geometry}
\usepackage{fontspec}
\setmainfont{Latin Modern Roman}
\usepackage{amsmath,amssymb,mathtools,unicode-math}
\setmathfont{Latin Modern Math}
\usepackage{enumitem,needspace,xurl,hyperref}
\setlength{\parindent}{0pt}
\setlength{\parskip}{5pt}
\setlength{\emergencystretch}{3em}
\newcommand{\N}{\mathbb{N}}
\newcommand{\Z}{\mathbb{Z}}
\newcommand{\Q}{\mathbb{Q}}
\newcommand{\R}{\mathbb{R}}
\newcommand{\C}{\mathbb{C}}
\newcommand{\F}{\mathbb{F}}
\newcommand{\A}{\mathbb{A}}
\newcommand{\PP}{\mathbb{P}}
\newcommand{\E}{\mathbb{E}}
\newcommand{\eps}{\varepsilon}
\newcommand{\dd}{\,\mathrm{d}}
\begin{document}
\section*{UnsolvedMath ID: 20001647}
\textbf{Source problem code:} catalogue entry
\subsection{Definitions and mathematical statement}
Let \(F_n\) be the free group of finite rank \(n\) and \(\operatorname{Out}(F_n)=\operatorname{Aut}(F_n)/\operatorname{Inn}(F_n)\). A group is large if a finite-index subgroup admits an epimorphism onto the nonabelian free group \(F_2\). Problem 3.5 asks whether \(\operatorname{Out}(F_n)\) is large. The extracted arrow \(\Gamma\to F_2\) must be read as surjective in the aid's intended question; an unrestricted homomorphism would make the trivial map an uninformative affirmative answer.

The aid reports a complete rank classification:
\[
\operatorname{Out}(F_n)\text{ is large}\quad\Longleftrightarrow\quad n=2\text{ or }n=3.
\]
For \(n\ge4\), it invokes modern Property (T) results for \(\operatorname{Aut}(F_n)\), then passage to \(\operatorname{Out}(F_n)\), to finite-index subgroups and to quotients. Since \(F_2\) does not have Property (T), no such epimorphism exists. Property (T) means that every unitary representation with almost invariant unit vectors has a nonzero invariant vector. The cases \(n=2,3\) use separate low-rank results. If rank one or zero is allowed, the finite groups \(\operatorname{Out}(F_1)\) and \(\operatorname{Out}(F_0)\) are not large. Full citations for the Property (T) and rank-three inputs are not supplied.

The aid gives two general transfer consequences. Let \(H\le G\) have finite index \(d\) and an epimorphism onto \(F_2\). For every integer \(m\ge1\), there is \(H_m\le H\) with
\[
[G:H_m]=dm,\qquad b_1(H_m;\Q)\ge m+1.
\]
The construction uses the preimage of an index-\(m\) free subgroup of \(F_2\), whose free rank is \(m+1\). These subgroups are not asserted to be normal in \(G\).

For every odd integer \(p\ge665\), the same hypothesis yields a normal subgroup \(N\trianglelefteq G\) such that \(G/N\) is infinite and has exponent dividing
\[
p\,\operatorname{lcm}(1,2,\ldots,d).
\]
Exponent dividing this integer means every element's corresponding power is the identity; it does not assert that all element orders equal the bound. One compatible description uses the core \(C=\bigcap_{g\in G}gHg^{-1}\) and its characteristic subgroup generated by \(p\)-th powers. The image \(G/C\le S_d\) has exponent dividing the displayed least common multiple, and the infinite-quotient input is the odd-exponent free Burnside theorem. The summary does not provide a complete authoritative citation for that input. The conclusion concerns a quotient of the ambient \(G\), not merely a torsion quotient of \(H\).

The source's surrounding discussion is retained as context: Thierry Coulbois asks which techniques apply; Bridson contrasts subgroup questions with train-track normal forms for elements; Juan Souto asks for an infinite torsion quotient and notes that this would obstruct bounded generation by finitely many cyclic subgroups, mentioning bounded cohomology. The transfer lemma supplies such quotients under the largeness hypothesis, but the negative largeness result for \(n\ge4\) does not by itself decide the separate infinite-torsion-quotient question in those ranks. The trailing comparison with hyperbolicity is discussion, not an additional precisely formulated theorem.

The registry labels the record solved and solved in literature; attempt 1 is a full solution, problem solved at the run, internal rigor proved and novelty confidence low. These labels concern the intended largeness question and supplied consequences, not automatic resolution of every surrounding discussion point. The aid is machine generated, not peer reviewed, and requires independent verification before citation.
\subsection{Short English statement}
Interpreting the arrow as an epimorphism, the aid reports that \(\operatorname{Out}(F_n)\) is large exactly in ranks two and three; Property (T) rules it out from rank four onward. It also gives quantitative Betti-number and infinite bounded-exponent quotient consequences for any large group, while preserving the surrounding source discussion as separate context.
\subsection{Sources}
\begin{itemize}
\item Local UnsolvedMath registry, permanent ID \(20001647\), source code \texttt{catalogue entry}; \texttt{aim-geometric-group-theory-notes.json}, \texttt{notes[155]}.
\item AIM, \emph{The geometry of the outer automorphism group of a free group}, Problem 3.5 and surrounding discussion involving Thierry Coulbois, Bridson and Juan Souto: \url{https://aimath.org/WWN/outerauto/outerauto.pdf}.
\item Modern Property (T) theorems for \(\operatorname{Aut}(F_n)\), low-rank largeness results and the odd-exponent Burnside input for every odd \(p\ge665\), as described in the supplied research summary. No full authors/publication/theorem references are provided; the original literature field is empty.
\item Supplied research attempt 1 and registry update August 14, 2026; no independent PDF inspection or current-literature verification is claimed.
\end{itemize}
\end{document}
